\begin{lemma}
\label{lemma-quasi-finite-finite-over-dense-open}
Let $f : Y \to X$ be a quasi-finite morphism.
There exists a dense open $U \subset X$ such that
$f|_{f^{-1}(U)} : f^{-1}(U) \to U$ is finite.
\end{lemma}

\begin{proof}
If $U_i \subset X$, $i \in I$ is a collection of opens such that the
restrictions $f|_{f^{-1}(U_i)} : f^{-1}(U_i) \to U_i$ are finite,
then with $U = \bigcup U_i$ the restriction $f|_{f^{-1}(U)} : f^{-1}(U) \to U$
is finite, see
Morphisms, Lemma \ref{morphisms-lemma-finite-local}.
Thus the problem is local on $X$ and we may assume that $X$ is affine.

\medskip\noindent
Assume $X$ is affine.
Write $Y = \bigcup_{j = 1, \ldots, m} V_j$ with $V_j$ affine.
This is possible since $f$ is quasi-finite and hence
in particular quasi-compact. Each $V_j \to X$ is quasi-finite
and separated. Let $\eta \in X$ be a generic point of an irreducible
component of $X$. We see from
Morphisms, Lemmas
\ref{morphisms-lemma-quasi-finite} and \ref{morphisms-lemma-generically-finite}
that there exists an open neighbourhood $\eta \in U_\eta$ such that
$f^{-1}(U_\eta) \cap V_j \to U_\eta$ is finite. We may choose $U_\eta$ such
that it works for each $j = 1, \ldots, m$.
Note that the collection of generic points of $X$ is dense in $X$.
Thus we see there exists a dense open $W = \bigcup_\eta U_\eta$
such that each $f^{-1}(W) \cap V_j \to W$ is finite.
It suffices to show that there exists a dense open $U \subset W$
such that $f|_{f^{-1}(U)} : f^{-1}(U) \to U$ is finite.
Thus we may replace $X$ by an affine open subscheme of $W$ and
assume that each $V_j \to X$ is finite.

\medskip\noindent
Assume $X$ is affine, $Y = \bigcup_{j = 1, \ldots, m} V_j$ with $V_j$ affine,
and the restrictions $f|_{V_j} : V_j \to X$ are finite.
Set
$$
\Delta_{ij} =
\Big(\overline{V_i \cap V_j} \setminus V_i \cap V_j\Big) \cap V_j.
$$
This is a nowhere dense closed subset of $V_j$ because it is the boundary
of the open subset $V_i \cap V_j$ in $V_j$. By
Morphisms, Lemma \ref{morphisms-lemma-image-nowhere-dense-finite}
the image $f(\Delta_{ij})$ is a nowhere dense closed subset of $X$. By
Topology, Lemma \ref{topology-lemma-nowhere-dense}
the union $T = \bigcup f(\Delta_{ij})$ is a nowhere dense closed
subset of $X$. Thus $U = X \setminus T$ is a dense open subset of $X$.
We claim that $f|_{f^{-1}(U)} : f^{-1}(U) \to U$ is finite.
To see this let $U' \subset U$ be an affine open.
Set $Y' = f^{-1}(U') = U' \times_X Y$,
$V_j' = Y' \cap V_j = U' \times_X V_j$. Consider the restriction
$$
f' = f|_{Y'} : Y' \longrightarrow U'
$$
of $f$. This morphism now has the property that
$Y' = \bigcup_{j = 1, \ldots, m} V'_j$ is an affine open covering,
each $V'_j \to U'$ is finite, and $V_i' \cap V_j'$ is (open and) closed
both in $V'_i$ and $V'_j$. Hence $V_i' \cap V_j'$ is affine, and the map
$$
\mathcal{O}(V'_i) \otimes_{\mathbf{Z}} \mathcal{O}(V'_j)
\longrightarrow
\mathcal{O}(V'_i \cap V'_j)
$$
is surjective. This implies that $Y'$ is separated, see
Schemes, Lemma \ref{schemes-lemma-characterize-separated}.
Finally, consider the commutative diagram
$$
\xymatrix{
\coprod_{j = 1, \ldots, m} V'_j \ar[rd] \ar[rr] & & Y' \ar[ld] \\
& U' &
}
$$
The south-east arrow is finite, hence proper, the horizontal arrow is
surjective, and the south-west arrow is separated. Hence by
Morphisms, Lemma \ref{morphisms-lemma-image-proper-is-proper}
we conclude that $Y' \to U'$ is proper. Since it is also quasi-finite,
we see that it is finite by Lemma \ref{lemma-characterize-finite},
and we win.
\end{proof}